\section{The Threshold of Christmas and the Mystery of Bethlehem}\label{sec:bethlehem}

Advent ends quietly in the afternoon of December~24. On that morning, the Church celebrates the final privileged weekday Mass of Advent (or the Fourth Sunday of Advent whenever December~24 falls on a Sunday), listening one last time to Zechariah's prophecy of ``the Orient from on high.'' Then, as twilight falls, Advent hands its watch over to the Solemnity of the Nativity of the Lord at Evening Prayer~I (\textit{Universal Norms on the Liturgical Year and the Calendar}, no.~40). The violet comes off, the white vestments and flowers fill the sanctuary, and the \textit{Gloria} held silent for four weeks breaks into song.

\subsection{The Four Masses of Christmas and the Proclamation}

To unfold a mystery too large for a single assembly, the postconciliar Roman Missal provides four distinct Christmas Masses, each with its own readings and prayers:

\begin{itemize}
  \item \textbf{The Vigil Mass} (late afternoon or early evening of December~24): Isaiah~62:1--5 promises that the Lord delights in Sion as a bridegroom rejoices in his bride; Psalm~89 (88) sings the Lord's mercies forever; Acts~13:16--17, 22--25 summarizes Israel's whole road to David's Son; and Matthew~1:1--25 gathers the genealogy and Joseph's obedience at the very threshold of the Birth.
  \item \textbf{Mass during the Night} (Midnight Mass): Isaiah~9:1--6 (Vulgate~9:2--7) proclaims ``The people that walked in darkness, have seen a great light \dots\ For a child is born to us, and a son is given to us''; Psalm~96 (95) calls heaven and earth to rejoice; Titus~2:11--14 declares that the grace of God has appeared to all men, training us to live soberly while awaiting the blessed hope; and Luke~2:1--14 narrates Caesar's decree, the manger, the shepherds, and the angels' \textit{Gloria}.
  \item \textbf{Mass at Dawn}: Isaiah~62:11--12 announces ``Behold thy Saviour cometh''; Psalm~97 (96) hails the light shining upon the just; Titus~3:4--7 remembers the kindness and love of God our Savior pouring out the bath of rebirth; and Luke~2:15--20 follows the shepherds hastening to Bethlehem to find Mary, Joseph, and the Infant lying in the manger.
  \item \textbf{Mass during the Day}: Isaiah~52:7--10 celebrates the feet of the herald bringing peace on the mountains; Psalm~98 (97) sings that all the ends of the earth have seen the salvation of our God; Hebrews~1:1--6 unveils the Son, the brightness of the Father's glory; and John~1:1--18 rises to the eternal summit of the feast: ``In the beginning was the Word \dots\ And the Word was made flesh and dwelt among us.''
\end{itemize}

Before Mass during the Night, many parishes sing the ancient \textbf{Proclamation of the Birth of Christ} (the \textit{Kalenda} from the \textit{Roman Martyrology}, printed in the 1908 Bute English rendering in Appendix~B). Moving through ages from the creation of the world, the Flood, Abraham, Moses, David, Daniel's weeks, the Olympiads, the founding of Rome, and the forty-second year of Caesar Augustus, the chant slows and rises as all kneel for the single sentence on which history hinges: ``Jesus Christ, Himself Eternal God and Son of the Eternal Father \dots\ was born of the Virgin Mary at Bethlehem of Juda made Man.'' Christ is not a myth floating outside time; He entered a measurable year under a real census to redeem every hour of human life.

\subsection{The House of Bread and the Manger}

The Church Fathers never treated Bethlehem's details as mere rustic scenery. Preaching his eighth \textit{Homily on the Gospels} in Rome on Christmas Day---apologizing for brevity because ``by God's mercy we are to say three Masses to-day''---Pope St.~Gregory the Great pointed straight to the census, the town's Hebrew name, and the animal feeding-trough:

\begin{quote}
Now the name Bethlehem signifieth ``the House of Bread,'' and thus it is the birth-place of Him Who hath said, ``I am the Living Bread, Which came down from heaven.'' We see then that this name of Bethlehem was prophetically given to the place where Christ was born, because it was there that He was to appear in the flesh by Whom the souls of the faithful are fed unto life eternal.
\end{quote}

Gregory presses the image even further (\textit{Homily~8.1}): because Scripture says ``All flesh is grass,'' the Word became flesh to turn our perishing grass into wheat, laying Himself in a manger (\textit{praesepium}) so that His faithful ones might be nourished with the grain of His Body and never starve for eternal food. Centuries later, St.~Bernard of Clairvaux (\textit{First Sermon on the Vigil of the Nativity}, I.1, 5--6) lingered over the same two words---``Jesus Christ, the Son of God, is born in Bethlehem of Juda''---remarking that ``Bethlehem is the house of bread; Juda signifies confession or praise,'' and adding a line about the poverty of the crib that overturns every worldly measure of wealth: ``There was an endless supply of these treasures in heaven, but poverty could not be found there. Earth abounded and superabounded in this kind of merchandise, and men knew not its value. The Son of God was desirous of it. He came down from heaven to make it His own, and so render it precious to us by His choice.''

Both St.~Ambrose of Milan (\textit{Exposition on Luke} II.41--42) and St.~Bede the Venerable (\textit{Exposition on Luke} 2:7) read every humbling circumstance of Luke~2:7 as an exact gift for humanity. Ambrose traces a series of deliberate exchanges---Christ became an infant so that we might grow into mature manhood, was wrapped in swaddling bands to loose us from the snares of death, lay in a manger so that we might stand at the altar, and had no other room in the inn so that we might inherit many mansions in heaven, calling that poverty his patrimony (\textit{meum ergo paupertas illa patrimonium est})---and sums up the mystery in words preserved in the 1843 Oxford \textit{Catena Aurea}:

\begin{quote}
I owe more, O Lord Jesus, to Thy sufferings that I am redeemed, than to Thy works that I am created. It were no advantage to be born, had it not advantaged me to be redeemed also. \dots\ On thy account then am I weak, in thee am I strong. On thy account am I poor, in thee am I rich. Consider not what thou seest, but acknowledge that thou art redeemed.
\end{quote}

\subsection{The Wonderful Exchange}

Beneath the tenderness of the manger lies the deepest doctrine of Catholic Christianity: the \textit{admirabile commercium}, the ``wonderful exchange'' sung on the Octave of Christmas (Appendix~B). God did not send a distant message of encouragement; He took our nature so that we might share His own life.

In the East, St.~Athanasius of Alexandria wrote the immortal summary in \textit{On the Incarnation of the Word} (\S54.3): ``For He was made man that we might be made God; and He manifested Himself by a body that we might receive the idea of the unseen Father; and He endured the insolence of men that we might inherit immortality.'' St.~Gregory of Nazianzus, preaching in Constantinople in his \textit{Oration~38 on the Theophany} (\S1, \S13, \S17), sang out in wonder: ``O new commingling; O strange conjunction; the Self-Existent comes into being, the Uncreate is created, That which cannot be contained is contained \dots\ He assumes the poverty of my flesh, that I may assume the richness of His Godhead.'' In his \textit{Third Theological Oration} (\textit{Oration~29.19}), Gregory traced the unbroken line from cradle to empty tomb: ``He was wrapped in swaddling clothes---but He took off the swathing bands of the grave by His rising again. He was laid in a manger---but He was glorified by Angels, and proclaimed by a star, and worshipped by the Magi.''

In the West, St.~Augustine echoed the same paradox in his Christmas preaching (\textit{Sermon~185.1}; \textit{Sermon~191.1}). Hailing the day on which Psalm~84:12 (\textit{Veritas de terra orta est}) is fulfilled, Augustine marvels in \textit{Sermon~185.1} that Truth, who abides in the Father's bosom and holds the world together, sprang from the earth to rest in a mother's bosom and be carried in a woman's arms. In \textit{Sermon~191.1}, he unfolds how the Word through whom all ages were made entered a single human day: the Maker of man became man (\textit{Homo factus, hominis factor}), so that the Ruler of the stars might nurse at a mother's breast, the Bread might hunger, the Fountain might thirst, and the Way might grow weary on the road. And Pope St.~Leo the Great, in his first Christmas sermon (\textit{Sermon~21.1, 3})---a homily still treasured across the Roman Office---turned that wonder straight into universal joy and moral courage:

\begin{quote}
Our Saviour, dearly-beloved, was born to-day: let us be glad. For there is no proper place for sadness, when we keep the birthday of the Life, which destroys the fear of mortality and brings to us the joy of promised eternity. \dots\ Let the saint exult in that he draws near to victory. Let the sinner be glad in that he is invited to pardon. Let the gentile take courage in that he is called to life. \dots\ Christian, acknowledge thy dignity, and becoming a partner in the Divine nature, refuse to return to the old baseness by degenerate conduct. Remember the Head and the Body of which thou art a member.
\end{quote}

\subsection{The Grace of Christmas Night in an Ordinary Home}

The mystery of the Incarnation never stays sealed inside ancient basilicas; it enters kitchen doors and bruised family hearts. In Chapter~V of her autobiography, \textit{The Story of a Soul}, St.~Th\'er\`ese of Lisieux recalls returning from Midnight Mass on Christmas~1886 as a sensitive thirteen-year-old who still wept easily at small slights. Climbing the stairs at home, she overheard her tired father sigh over the childish custom of leaving her shoes by the fireplace. In any other year she would have burst into tears and ruined the evening. Instead, on that night, the Infant Jesus who had made Himself weak and small out of love poured sudden strength into her soul: she mastered her tears, went downstairs with a radiant smile, and opened her gifts in quiet self-forgetfulness. Th\'er\`ese looked back on Christmas~1886 as the night of her decisive conversion, the hour when charity entered her heart and launched her upon the apostolic race to win souls for Christ.

That is what Advent has been training us for from the first evening of late November. Every quiet morning with Isaiah, every Confession, every rose candle lit in winter dusk, every O~Antiphon whispered at table, and every hidden act of mercy toward the poor has had a single end: that when the Word made flesh arrives on Christmas Night, He may find in us not merely a decorated parlor, but a wakeful, welcoming Bethlehem---a house of bread, confession, and praise.
